\documentclass[12pt,a4paper]{article}

% --------------------------------------------------
% Packages
% --------------------------------------------------
\usepackage[utf8]{inputenc}
\usepackage[T1]{fontenc}
\usepackage{amsmath,amssymb,amsthm}
\usepackage{mathtools}
\usepackage{geometry}
\usepackage{booktabs}
\usepackage{graphicx}
\usepackage{hyperref}

\geometry{margin=1in}

% --------------------------------------------------
% Theorem environments
% --------------------------------------------------
\newtheorem{theorem}{Theorem}
\newtheorem{lemma}{Lemma}
\newtheorem{proposition}{Proposition}
\newtheorem{corollary}{Corollary}

\theoremstyle{definition}
\newtheorem{definition}{Definition}
\newtheorem{example}{Example}

% --------------------------------------------------
% Document information
% --------------------------------------------------
\title{Mathematical Typesetting Demonstration}
\author{Dr. Bosede O. Fagbemigun}
\date{\today}

\begin{document}

\maketitle

\begin{abstract}
This document demonstrates selected LaTeX techniques commonly
used in mathematical research and academic writing, including
mathematical notation, theorem environments, proofs, equations,
matrices, tables, and cross-referencing.
\end{abstract}

\section{Introduction}

LaTeX provides a structured environment for preparing mathematical
and scientific documents. It is particularly useful when a document
contains complex mathematical notation, numbered equations,
theorem environments, tables, figures, and bibliographic references.

\section{Mathematical Expressions}

Consider the following inequality:

\begin{equation}
\left|
\frac{1}{b-a}\int_a^b f(x)\,dx
\right|
\leq
\frac{1}{b-a}\int_a^b |f(x)|\,dx.
\label{eq:basic-inequality}
\end{equation}

Equation~\eqref{eq:basic-inequality} illustrates the use of numbered
equations and internal cross-referencing.

A more structured expression can be written as

\begin{align}
F(x)
&= x^2 + 2x + 1 \\
&= (x+1)^2.
\label{eq:factorization}
\end{align}

The `align` environment allows multiple lines of mathematics to be
aligned clearly.

\section{Theorem and Proof}

\begin{definition}
A function $f:[a,b]\rightarrow\mathbb{R}$ is said to be convex if
\[
f(tx+(1-t)y)
\leq
tf(x)+(1-t)f(y)
\]
for all $x,y\in[a,b]$ and $t\in[0,1]$.
\end{definition}

\begin{theorem}
Let $f:[a,b]\rightarrow\mathbb{R}$ be convex. Then for every
$x,y\in[a,b]$ and $t\in[0,1]$,
\[
f(tx+(1-t)y)
\leq
tf(x)+(1-t)f(y).
\]
\end{theorem}

\begin{proof}
The result follows directly from the definition of convexity.
The definition states precisely that the value of the function at
a convex combination of two points does not exceed the corresponding
convex combination of the function values.
\end{proof}

\section{Matrices}

A matrix can be typeset using the \texttt{pmatrix} environment:

\[
A =
\begin{pmatrix}
a_{11} & a_{12} & a_{13}\\
a_{21} & a_{22} & a_{23}\\
a_{31} & a_{32} & a_{33}
\end{pmatrix}.
\]

Its determinant may be represented by

\[
\det(A)
=
\begin{vmatrix}
a_{11} & a_{12}\\
a_{21} & a_{22}
\end{vmatrix}.
\]

\section{Tables}

Table~\ref{tab:example} demonstrates a simple academic table.

\begin{table}[htbp]
\centering
\caption{Illustrative mathematical quantities}
\label{tab:example}
\begin{tabular}{ccc}
\toprule
Quantity & Symbol & Value \\
\midrule
Length & $L$ & $10$ \\
Width & $W$ & $5$ \\
Area & $A$ & $50$ \\
\bottomrule
\end{tabular}
\end{table}

\section{Cross-References}

LaTeX makes it possible to reference mathematical objects
automatically. For example, Equation~\eqref{eq:factorization}
shows a factorization, while Table~\ref{tab:example} presents
illustrative numerical quantities.

\section{Conclusion}

This demonstration illustrates several features commonly required
in mathematical and scientific manuscripts:

\begin{itemize}
    \item mathematical notation and equations;
    \item theorem, definition, and proof environments;
    \item matrices and structured mathematical expressions;
    \item tables and captions;
    \item labels and cross-references;
    \item structured academic document preparation.
\end{itemize}

The source file can be compiled using a standard LaTeX distribution
such as MiKTeX or through an online collaborative environment such
as Overleaf.

\end{document}
